\documentclass[11pt]{article}
% teacher_preamble.tex — shared preamble for the chalk-teaching TEACHER'S NOTES.
% These are standalone, instructor-facing lesson plans (article class), separate
% from the student lecture notes.  Compile each teaching/lecNN_teacher_notes.tex
% directly with pdfLaTeX.
\usepackage[T1]{fontenc}
\usepackage[utf8]{inputenc}
\usepackage[margin=1in]{geometry}
\usepackage{amsmath,amssymb,amsthm,mathtools}
\usepackage{braket}
\usepackage{bm}
\usepackage{graphicx}
\usepackage{enumitem}
\usepackage{xcolor}
\usepackage{tikz}
\usetikzlibrary{arrows.meta}
\usepackage{quantikz}
\usepackage[most]{tcolorbox}
\usepackage[colorlinks=true,linkcolor=blue,urlcolor=blue]{hyperref}

% ---- Notation macros (consistent with the course preamble.tex) ----
\newcommand{\R}{\mathbb{R}}
\newcommand{\C}{\mathbb{C}}
\newcommand{\N}{\mathbb{N}}
\newcommand{\Z}{\mathbb{Z}}
\newcommand{\E}{\mathbb{E}}
\newcommand{\Prob}{\mathbb{P}}
\newcommand{\Hilb}{\mathcal{H}}
\newcommand{\ii}{\mathrm{i}}
\newcommand{\dd}{\mathrm{d}}
\newcommand{\ee}{\mathrm{e}}
\DeclareMathOperator{\tr}{tr}
\DeclareMathOperator{\Tr}{Tr}
\DeclareMathOperator{\Var}{Var}
\DeclareMathOperator{\diag}{diag}
\newcommand{\norm}[1]{\left\lVert #1 \right\rVert}
\newcommand{\abs}[1]{\left\lvert #1 \right\rvert}
\newcommand{\bigo}{\mathcal{O}}
\newcommand{\id}{\mathbb{I}}
\newcommand{\dg}{^{\dagger}}
% quantikz also defines \proj (a circuit element we do not use); override it with our projector.
\let\proj\relax
\newcommand{\proj}[1]{\ket{#1}\!\bra{#1}}
\newcommand{\Xgate}{\mathsf{X}}
\newcommand{\Ygate}{\mathsf{Y}}
\newcommand{\Zgate}{\mathsf{Z}}
\newcommand{\Hgate}{\mathsf{H}}
\newcommand{\CNOT}{\mathrm{CNOT}}
\newcommand{\Uf}{U_f}

% ---- Theorem-like environments (numbered per section, article style) ----
\theoremstyle{definition}
\newtheorem{definition}{Definition}[section]
\newtheorem{example}[definition]{Example}
\theoremstyle{plain}
\newtheorem{theorem}[definition]{Theorem}
\newtheorem{lemma}[definition]{Lemma}
\newtheorem{proposition}[definition]{Proposition}
\theoremstyle{remark}
\newtheorem{remark}[definition]{Remark}

% ---- Teaching-specific structure ----
% Session block header with a time budget, e.g. \block{Qubits and the Bloch sphere}{45}
\newcommand{\block}[2]{%
  \bigskip\noindent
  \begin{tcolorbox}[colback=blue!6,colframe=blue!55!black,boxrule=0.6pt,arc=2pt,
    left=8pt,right=8pt,top=5pt,bottom=5pt]
  \textbf{\large #1}\hfill\textbf{[\,#2\,min]}
  \end{tcolorbox}}

% Board-plan cue: what to write/draw on the board.
\newtcolorbox{board}{colback=gray!8,colframe=gray!55!black,title={On the board},
  fonttitle=\bfseries\small,boxrule=0.5pt,arc=2pt,left=6pt,right=6pt,top=3pt,bottom=3pt}

% Graphic-demonstration cue (what to sketch).
\newcommand{\sketch}[1]{\par\smallskip\noindent{\color{teal}\textbf{Sketch:}} #1\par\smallskip}

% Optional / advanced material that may be skipped (but indicated).
\newtcolorbox{advanced}{colback=orange!7,colframe=orange!65!black,
  title={Advanced --- may be skipped (mention only)},
  fonttitle=\bfseries\small,boxrule=0.5pt,arc=2pt,left=6pt,right=6pt,top=3pt,bottom=3pt}

% Timing hint inline, e.g. \tmark{10} for a ~10-minute chunk.
\newcommand{\tmark}[1]{{\small\color{blue!55!black}(\,$\approx$#1\,min)}}

% In-class problem (numbered per section) and a teacher-facing answer.
\newtheorem{inclass}[definition]{In-class problem}
\newenvironment{answer}{\par\smallskip\noindent\textit{Answer.}\ }{\par\smallskip}

\providecommand{\QFT}{\mathrm{QFT}}
\title{\textbf{Lecture 5 --- Applications of Phase Estimation}\\[0.3em]\large Teacher's Notes (chalk, one 3$\times$45\,min meeting)}
\author{Course Instructor}
\date{}
\begin{document}
\maketitle

\noindent\textbf{How to use these notes.} A chalk plan for one three-hour meeting in
three 45-minute blocks. \textsf{On the board} panels give what to write;
\textcolor{teal}{\textbf{Sketch}} cues what to draw; \tmark{n} paces each chunk. The
spine is a single idea---\emph{QPE reads an eigenphase}---applied three ways; keep
pointing back to Lecture~4. In-class problems have teacher answers; orange
\emph{Advanced} boxes are optional.

\medskip
\noindent\textbf{Learning goals (by the end a student can):}
\begin{itemize}[leftmargin=1.4em,nosep]
  \item Estimate a ground-state energy by running QPE on $\ee^{-\ii Ht}$, and explain why the cost is set by the input/ground-state overlap.
  \item Recognize amplitude estimation as QPE of the Grover iterate, and state its $\bigo(1/\epsilon)$ vs.\ $\bigo(1/\epsilon^2)$ quadratic advantage.
  \item Cast order finding as phase estimation of modular multiplication and recover the order by continued fractions (the quantum core of Shor).
  \item State the Heisenberg limit and the no-fast-forwarding caveat at the level of their consequences.
\end{itemize}

\medskip
\noindent\textbf{Prerequisites.} Lecture~4 (QPE, inverse QFT, precision bound);
Lecture~3 (Grover iterate); basic number theory ($\gcd$, modular arithmetic, continued
fractions).

\medskip
\noindent\textbf{Materials.} Cheat strip: $H\ket{\psi_j}=\lambda_j\ket{\psi_j}$,
ground energy $\lambda_0$, gap $\Delta=\lambda_1-\lambda_0$, overlap $p_0=|\braket{\phi|\psi_0}|^2$;
colored chalk for ``signal'' (true value) vs.\ ``estimate.''

% =====================================================================
\block{Block 1 --- Ground-state energy estimation}{45}
% =====================================================================
\noindent\emph{Objective:} eigenphases of $\ee^{-\ii Ht}$ are energies; QPE on a trial
state samples energy $\lambda_j$ with probability $|\braket{\phi|\psi_j}|^2$.

\medskip
\noindent\textbf{Pacing:} eigenphase $\to$ energy \tmark{12}; the overlap governs cost
\tmark{12}; min-over-runs estimator \tmark{10}; in-class \tmark{3}; buffer \tmark{8}.

\subsection*{1. From eigenphase to energy \tmark{12}}
\begin{board}
Run QPE with $U=\ee^{-\ii Ht}$. Its eigenphases are $\varphi_j=-\lambda_j t/2\pi \pmod1$,
so a $t$-ancilla QPE on $\ket{\psi_j}$ returns $\widetilde\varphi_j\approx\varphi_j$ and
hence $\widetilde\lambda_j=-2\pi\widetilde\varphi_j/t$. Choose the evolution time so all
$\lambda_j t$ fit in $[0,2\pi)$ (rescale $H$) to avoid aliasing.
\end{board}
Say aloud: ``energy = phase accumulated per unit time''; QPE is a quantum spectrometer.
Implementing $\ee^{-\ii Ht}$ (Hamiltonian simulation) is Part~III---here it is a black box.

\subsection*{2. The overlap governs the cost \tmark{12}}
\begin{board}
On a trial state $\ket\phi=\sum_j c_j\ket{\psi_j}$, QPE returns $\lambda_j$ with
probability $|c_j|^2$ and collapses onto $\ket{\psi_j}$. To hit the \emph{ground} energy
$\lambda_0$ you need outcome $j=0$: expected repetitions $\sim 1/p_0$, $p_0=|c_0|^2$.
A good classical/variational guess for $\ket\phi$ (large $p_0$) is essential.
\end{board}
\sketch{Spectrum as a comb of lines $\lambda_0<\lambda_1<\cdots$; the trial state's bars
$|c_j|^2$ over them; QPE ``samples a line'' with those weights.}

\subsection*{3. The min-over-runs estimator \tmark{10}}
\begin{board}
Repeat QPE $R$ times; each run yields some $\widetilde\lambda_{j}$. Report the
\emph{minimum}. With $p_0$ not tiny and $R=\bigo(\tfrac1{p_0}\log\tfrac1\delta)$ runs, the
smallest observed estimate equals $\widetilde\lambda_0$ (to QPE precision) except with
probability $\delta$---the ground state is the lowest line, so its sample is the min.
\end{board}

\begin{inclass}[Repetitions for the ground energy]\label{ic:l5-overlap}
A trial state has ground-state overlap $p_0=0.04$. Roughly how many QPE runs are needed
so that at least one hits $\lambda_0$ with probability $\ge0.99$?
\end{inclass}
\begin{answer}
$\Prob(\text{miss all }R)=(1-p_0)^R\le0.01\Rightarrow R\ge\ln(0.01)/\ln(0.96)\approx113$.
So $\sim\!10^2$ runs; cost scales as $1/p_0$ (QSVT filtering in \S Block~3 improves this to $1/\sqrt{p_0}$).
\end{answer}

% =====================================================================
\block{Block 2 --- Amplitude estimation}{45}
% =====================================================================
\noindent\emph{Objective:} the ``good-subspace weight'' is an eigenphase of the Grover
iterate, so QPE estimates it with a quadratic speedup over sampling.

\medskip
\noindent\textbf{Pacing:} setup \& the Grover-iterate eigenphase \tmark{16}; query count
\& quadratic advantage \tmark{14}; applications \tmark{9}; in-class \tmark{6}.

\subsection*{1. The weight as an eigenphase \tmark{16}}
\begin{board}
Let $\mathcal A\ket0=\ket\Psi=\sin\tfrac\theta2\ket{\text{good}}+\cos\tfrac\theta2\ket{\text{bad}}$,
so the good-subspace probability is $a=\sin^2\tfrac\theta2$. The Grover iterate
$Q=\mathcal A S_0\mathcal A^\dagger S_{\text{good}}$ rotates by $\theta$ in the 2-D
good/bad plane; its eigenvalues are $\ee^{\pm\ii\theta}$. \textbf{Amplitude estimation}
$=$ QPE on $Q$ to read $\theta$, hence $a=\sin^2\tfrac\theta2$.
\end{board}
Point back: this is exactly the Grover geometry of Lecture~3, now \emph{measured} rather
than iterated to a target.

\subsection*{2. Query count and the quadratic advantage \tmark{14}}
\begin{board}
A $t$-bit QPE on $Q$ uses $\bigo(2^t)$ applications of $Q$ and resolves $\theta$ to
$\bigo(2^{-t})$, so estimating $a$ to additive error $\epsilon$ costs
\[
  \bigo(1/\epsilon)\ \text{queries}\quad\text{vs.}\quad\bigo(1/\epsilon^2)\ \text{for Monte Carlo.}
\]
A \emph{quadratic} speedup in precision---the same $\sqrt{\ }$ that powers Grover.
\end{board}

\subsection*{3. Applications \tmark{9}}
\begin{board}
\emph{Quantum counting}: $\mathcal A=H^{\otimes n}$, good $=$ marked set; $a=M/N$ gives the
number of solutions $M$---and the right Grover iteration count without knowing $M$.
\emph{Monte Carlo}: estimating a mean $\E[f]$ to error $\epsilon$ in $\bigo(1/\epsilon)$
queries (quadratically fewer samples).
\end{board}

\begin{inclass}[Counting speedup]\label{ic:l5-ae}
To estimate a probability $a$ to additive error $\epsilon=10^{-3}$: how many samples does
classical Monte Carlo need (order of magnitude), and how many Grover-iterate queries does
amplitude estimation need?
\end{inclass}
\begin{answer}
Monte Carlo $\sim1/\epsilon^2=10^{6}$ samples; amplitude estimation $\sim1/\epsilon=10^{3}$
queries to $Q$---a thousand-fold reduction here, the quadratic advantage.
\end{answer}

% =====================================================================
\block{Block 3 --- Order finding, Shor, and limits}{45}
% =====================================================================
\noindent\emph{Objective:} factoring reduces to order finding; order finding is QPE of a
modular-multiplication unitary; continued fractions recover the order.

\medskip
\noindent\textbf{Pacing:} factoring $\to$ order finding \tmark{10}; order finding as QPE
\tmark{14}; continued fractions \& speedup \tmark{12}; limits (brief) \tmark{6}; in-class \tmark{3}.

\subsection*{1. Factoring reduces to order finding \tmark{10}}
\begin{board}
To factor $M$: pick $1<a<M$ with $\gcd(a,M)=1$; let $r=\mathrm{ord}_M(a)$ be the least
$r>0$ with $a^r\equiv1\pmod M$. If $r$ is even and $a^{r/2}\not\equiv-1$, then
$\gcd(a^{r/2}\pm1,\,M)$ are nontrivial factors. The hard part is finding $r$---do it quantumly.
\end{board}

\subsection*{2. Order finding as phase estimation \tmark{14}}
\begin{board}
Modular-multiplication unitary $U_a\ket{y}=\ket{ay\bmod M}$. Its eigenstates
\[
  \ket{u_s}=\tfrac1{\sqrt r}\sum_{k=0}^{r-1}\ee^{-\ii2\pi sk/r}\ket{a^k\bmod1}
  \quad\text{have eigenphase } \varphi_s=\tfrac{s}{r}.
\]
Starting the eigenstate register in $\ket1=\tfrac1{\sqrt r}\sum_s\ket{u_s}$, QPE returns a
random $s/r$ (each $s$ equally likely). Controlled-$U_a^{2^j}$ is modular exponentiation.
\end{board}
Say aloud: we never prepare a single $\ket{u_s}$---$\ket1$ is their uniform superposition,
so one QPE run yields $\widetilde{s/r}$ for a random $s$.

\subsection*{3. Recovering $r$ and the speedup \tmark{12}}
\begin{board}
From the estimate $\widetilde{s/r}$ (to $2n+1$ bits, $n=\lceil\log_2 M\rceil$), the
\emph{continued-fraction} expansion returns the fraction $s/r$ in lowest terms, hence $r$
(when $\gcd(s,r)=1$; retry otherwise---$\bigo(1)$ expected repeats). Cost: $\bigo(n^2)$
for the QFT $+$ modular exponentiation, i.e.\ $\mathrm{poly}(n)$ --- \emph{exponentially}
faster than the best known classical factoring.
\end{board}

\begin{inclass}[Continued fractions]\label{ic:l5-cf}
QPE on $U_a$ with $M=15$, $a=7$ returns the estimate $\widetilde{s/r}=0.25$. What are the
candidate $s/r$, the order $r$, and the resulting factors of $15$?
\end{inclass}
\begin{answer}
$0.25=\tfrac14\Rightarrow s/r=1/4$, so $r=4$. Check $7^4=2401\equiv1\pmod{15}$. Then
$a^{r/2}=7^2=49\equiv4$, and $\gcd(4\pm1,15)=\gcd(5,15),\gcd(3,15)=5,3$: factors $15=3\times5$.
\end{answer}

\subsection*{4. Limits --- state them, don't prove \tmark{6}}
\begin{board}
\emph{Heisenberg limit}: $N_U$ uses of $U$ buy phase precision $\sim1/N_U$ (optimal).
\emph{Ground-state prep} (QSVT spectral filter, Lecture~9): cost $\sim1/\sqrt{p_0}$ and
$\sim1/\Delta$---a Grover-type $\Omega(1/\sqrt{p_0})$, $\Omega(1/\Delta)$ lower bound shows
this is near-optimal. \emph{No fast-forwarding}: generic $\ee^{-\ii Ht}$ needs
$\Omega(t)$ work---recurs in Part~III.
\end{board}

\begin{advanced}
\textbf{Eigenstate preparation by filtering.} A polynomial $p(H)\approx$ projector onto
$\lambda\le\lambda_0+\Delta/2$, applied by QSVT (Lecture~9), amplifies the ground
component from amplitude $\sqrt{p_0}$ to $\bigo(1)$ using $\bigo(1/\sqrt{p_0})$ calls and
degree $\bigo(1/\Delta)$---quadratically better in $p_0$ than the $1/p_0$ sampling of
Block~1. Mention only; the machinery is QSVT.
\end{advanced}

\end{document}
